\documentclass[10pt,a4paper]{amsart}
\usepackage[margin=19mm]{geometry}
\usepackage{amsmath,amssymb,mathtools}
\usepackage[scr=rsfs,cal=euler]{mathalfa}
\usepackage{xfrac}
\usepackage[unicode,hidelinks,bookmarks=false]{hyperref}
\newcommand{\ie}{i.e.\ }
\newcommand{\cf}{cf.\ }
\newcommand{\resp}{respectively\ }
\newcommand{\Subsection}{\subsection}
\providecommand{\nobreakdash}{\nobreak\hbox{-}}
\setlength{\emergencystretch}{2em}
\multlinegap0pt
\usepackage{tikz}
\usetikzlibrary{cd}
\usepackage{amssymb}
\usepackage{mathtools}
\usepackage{tikz}
\newtheorem{prop}{Proposition}
\newtheorem{ex}{Example}
\newtheorem{defn}{Definition}
\newtheorem{thm}{Theorem}
\newcommand{\CC}{\mathbb{C}}
\newcommand{\GL}[1]{GL_{#1}(\mathbb{F}_q)}
\def\IM#1{\textcolor[rgb]{0.00,0.00,0.75}{[I: #1]}}
\DeclareMathOperator{\R}{\mathcal{R}}
\title{Radon transform for $GL_n(\mathbb{F}_q)$}
\author{Ivan Motorin, Kai Yamashita}
\date{Source: arXiv:2606.18002v2 (CC BY 4.0)}
\begin{document}
\maketitle

\noindent\emph{Source and licence.} Radon transform for $GL_n(\mathbb{F}_q)$. Version arXiv:2606.18002v2 (CC BY 4.0), distributed by arXiv under the Creative Commons Attribution 4.0 International licence, http://creativecommons.org/licenses/by/4.0/, as recorded in the arXiv licence metadata for this exact version and shown on the abstract page https://arxiv.org/abs/2606.18002v2. Full licence notice: \texttt{licenses/arxiv-2606.18002-CC-BY-4.0.txt}.\par\smallskip

\noindent\textit{Editorial scope.} Counted unit: the article's thm eigenvalues\_gl2 (source0.tex lines 741--743). The article's complete original proof of it is delivered with it (source0.tex lines 745--1042, one \texttt{proof} environment of 298 lines). No mathematics is written, repaired or re-derived: the statement and the proof are the article's own lines, in the article's own notation, and no retained line is substituted or re-flowed.  Retained uncounted as the source-local context the proof needs: lines 230--236; lines 282--289; lines 511--524; lines 534--548.  Nothing was omitted from any retained span, so the package carries no omission record.  No retained line contains a citation, so the wrapper carries no bibliography and no bibliography entry.  The retained lines contain the article's own diagram code, which the wrapper typesets with the standard tikz packages; no image file is delivered and the package declares no asset dependency.  Editorial formatting only, in addition: an AMS \texttt{amsart} preamble that restates the article's own declarations and macros used by the retained lines, namely the environments \texttt{prop}, \texttt{ex}, \texttt{defn}, \texttt{thm} on separate counters, together with the standard packages those lines need.  The retained environments print in this document as Proposition 1, Example 1, Definition 1, Theorem 1 in order of appearance, whereas the article prints the same items with the numbering of its own full document.  The complete native source of the article as distributed in the arXiv e-print for this exact version is bundled unchanged as \texttt{sources/203056/source0.tex}.
\begin{prop}\label{Functions-group_algebra_correspondence}
    Let us denote by $\phi$ the following map
    \begin{equation}\label{fun_iso}
        \phi: Fun(G,k) \rightarrow k[G], \quad f \mapsto \sum_{g\in G} f(g) g.
    \end{equation}
    Then $\phi$ is an isomorphism of algebras $(Fun(G,k),m')$ and $(k[G],m)$. Furthermore, the algebra isomorphism $\phi$ is an isomorphism of left $k[G \times G]$-modules.
\end{prop}

\begin{ex}\label{subgroup-projector}
    One example of an idempotent corresponds to a subgroup $H$ of a group $G$, defined in the group algebra of $G$ as
    
    \begin{equation*}
        e_H:= \frac{1}{|H|}\sum_{h\in H}h.
    \end{equation*}
    Here we assume that $|H|$ is invertible in $k$.
\end{ex}

\begin{defn}[{\bf Radon transform}]\label{Radon_transform}
    Consider the following diagram. %\IM{\\
    %https://tikzcd.yichuanshen.de for diagrams}
    \begin{equation*}
    %\shorthandoff{"}%
        \begin{tikzcd}
            & {Fun(GL_n(\mathbb{F}_q),\mathbb{C})} \arrow[rd, "\pi_*"', shift right] \arrow[ld, "\pi^{op}_*"', shift right] & \\
            {Fun(GL_n(\mathbb{F}_q)/U_{\lambda}^{op},\mathbb{C})} \arrow[ru, "\pi^{op *}"', shift right] &                                                                                                               & {Fun(GL_n(\mathbb{F}_q)/U_{\lambda},\mathbb{C})} \arrow[lu, "\pi^*"', shift right]
        \end{tikzcd}
.
%\shorthandon{"}%
    \end{equation*}
    Where the maps $\pi_*,\pi^*,\pi_*^{op},\pi^{op*}$ are given by the "pull back" and "push forward" operators. For instance, $\pi^*$ is the operator that is defined by $\pi^*(f)(g) = f(gU_\lambda)$, and $\pi_*$ is the operator $\pi_*(f)(gU_\lambda) = \sum_{u \in U_\lambda} f(gu)$. Using this notation, we denote $R = \pi^{*} \pi^{op}_*$ and $R' = \pi^{op*} \pi_*$. We will study the Radon operator which we will denote as $\R:=R' \circ R$.
\end{defn}

\begin{prop}\label{Radon Expansion}
    Under the map $\phi$ from the Proposition \ref{Functions-group_algebra_correspondence} the Radon transform $\R$ from Definition \ref{Radon_transform} corresponds to the endomorphism of the subspace $\CC[\GL{n}/U_{\lambda}]$ of right $U_{\lambda}$-invariants in $\CC[\GL{n}]$. Namely, it is given by the operation of the right multiplication
    \begin{equation}
        R_{\R}: \CC[\GL{n}]e_{U_{\lambda}} \rightarrow \CC[\GL{n}]e_{U_{\lambda}}, \quad g\cdot e_{U_{\lambda}} \mapsto g\cdot |U_{\lambda}|^2e_{U_{\lambda}}e_{U_{\lambda}^{op}}e_{U_{\lambda}}.
    \end{equation}
    From now on we work with the group element $\R:=|U_{\lambda}|^2e_{U_{\lambda}}e_{U_{\lambda}^{op}}e_{U_{\lambda}} \in \CC[\GL{n}]$.
    
    % \IM{Try to use the notation above} We use the notation $V_\lambda$ and $V_\lambda^{op}$ to refer to $\sum_{u\in U_\lambda} u$ and $\sum_{u^{op} \in U_\lambda^{op}} u^{op}$ respectively, so the Radon transform is equal to $g V_\lambda V_\lambda^{op} U_\lambda$ .\\

    %Since Radon transform projects the group algebra $\CC[\GL{n}]$ onto the subspace of %Sometimes, it is easier to view everything in the group algebra $\CC[\GL{n}]$
    %with right $U_\lambda$-invariant elements of $\mathbb{C}[\GL{n}]$ (i.e. we apply $\pi^*$ and $\pi^{op*}$ as needed). In this case, the Radon transform is equal to 
    %\begin{equation}
    %    R_{\R}|_{\CC[\GL{n}]e_{U_{\lambda}}}: g\cdot e_{U_{\lambda}} \mapsto |U_{\lambda}|^2 g \cdot e_{U_{\lambda}}e_{U_{\lambda}^{op}}e_{U_{\lambda}}.
    %\end{equation}
\end{prop}

\begin{thm}\label{eigenvalues_gl2}
    The only possible eigenvalues of the Radon transform on $\GL{2}$ are $1$, $q$, and $q^2$.
\end{thm}

\begin{proof}
Let $\lambda=(1,1)$ and $e_{U_{\lambda}}=\frac{1}{|U_{\lambda}|}\sum_{u \in U_\lambda} u$ as defined in \ref{subgroup-projector}. Recall from the Definition \ref{Radon Expansion} that $\R =q^2 e_{U_{\lambda}}e_{U_{\lambda}^{op}}e_{U_{\lambda}}$.\\

%Specifically, $\R (g) = \frac{1}{|U_\lambda|} g V_\lambda V^{op}_\lambda V_\lambda$.\\

We explicitly write out the minimal polynomial of $\R$. For this we will find similarities between the elements $q^3 e_{U_{\lambda}}e_{U_{\lambda}^{op}}e_{U_{\lambda}}$ and $q^3 e_{U_{\lambda}^{op}}e_{U_{\lambda}}e_{U_{\lambda}^{op}}$. We have

\[q^3 e_{U_{\lambda}}e_{U_{\lambda}^{op}}e_{U_{\lambda}} = \sum_{a, b, c \in \mathbb{F}_q} 
\begin{pmatrix}
1 & a \\
0 & 1 
\end{pmatrix}
\begin{pmatrix}
1 & 0 \\
b & 1 
\end{pmatrix}
\begin{pmatrix}
1 & c \\
0 & 1 
\end{pmatrix}
= \sum_{a, b, c \in \mathbb{F}_q} 
\begin{pmatrix}
1+ab & c+a+abc \\
b & bc+1 
\end{pmatrix}.\]

Similarly,

\[q^3 e_{U_{\lambda}^{op}}e_{U_{\lambda}}e_{U_{\lambda}^{op}} = \sum_{a', b', c' \in \mathbb{F}_q} 
\begin{pmatrix}
1 & 0 \\
a' & 1 
\end{pmatrix}
\begin{pmatrix}
1 & b' \\
0 & 1 
\end{pmatrix}
\begin{pmatrix}
1 & 0 \\
c' & 1 
\end{pmatrix}
= \sum_{a', b', c' \in \mathbb{F}_q} 
\begin{pmatrix}
1+b'c' & b' \\
c'+a'+a'b'c' & a'b'+1 
\end{pmatrix}.\]

Setting both matrices to be equal, we get the system of equations:
\[
    b'c' = ab, \quad a'b' = bc, \quad b' = a+c+abc, \quad b = a'+c'+a'b'c'.
\]
%\begin{align*}
%b'c' &= ab\\
%a'b' &= bc\\
%b' &= a+c+abc\\
%b &= a'+c'+a'b'c'
%\end{align*}

Using the first three equations, we get the following solutions
\[
    a' = \frac{bc}{a+c+abc}, \quad b' = a+c+abc, \quad c' = \frac{ab}{a+c+abc}.
\]
%\begin{align*}
%a' = \frac{bc}{a+c+abc}\\
%b' = a+c+abc\\
%c' = \frac{ab}{a+c+abc}
%\end{align*}
Here, solutions exist only if $a+c+abc \ne 0$. Note that this corresponds to the matrices where $b' = 0$. This makes sense as if $b' = 0$, then the corresponding summands of $\R$ are lying exclusively in $U^{op}_\lambda$. There are multiple choices for $a'$ and $c'$ for every element, so it makes sense that we cannot find a single solution for $a'$ $b'$ and $c'$. In addition, we also need to account for the matrices where $b=0$, or when $c'+a'+a'b'c' = 0$. This is again a problem as there are multiple values of $a$ and $c$ that satisfy this equality, so we need to remove these matrices before comparing the two expressions.\\

The matrices where $b' = 0$ are equivalent to the matrices where $a+c+abc = 0$. This can be split into two cases, one where $a$ and $c$ are both $0$, and the other where $b = -\frac{a+c}{ac}$.\\

The other case where $b = 0$ is equivalent. Then, we also have to make sure we do not double-count the cases where both $b = 0$ and $b' = 0$, which is simply just the cases where the matrix is the identity.\\

%Having been careful about what we subtract, we get that the following
From the observations above we get that

\[q^3 e_{U_{\lambda}}e_{U_{\lambda}^{op}}e_{U_{\lambda}} - \sum_{a, c \in \mathbb{F}_q} 
\begin{pmatrix}
1 & a+c \\
0 & 1 
\end{pmatrix} - \sum_{b \in \mathbb{F}_q}
\begin{pmatrix}
1 & 0 \\
b & 1 
\end{pmatrix}
- \sum_{a, c \in \mathbb{F}_q^{\times}}
\begin{pmatrix}
-a/c & 0 \\
-(a+c)/ac & -c/a 
\end{pmatrix}
+\sum_{a \in \mathbb{F}_q}
\begin{pmatrix}
1 & 0 \\
0 & 1 
\end{pmatrix}\]

is equal to

\[q^3 e_{U_{\lambda}^{op}}e_{U_{\lambda}}e_{U_{\lambda}^{op}} - \sum_{a', c' \in \mathbb{F}_q} 
\begin{pmatrix}
1 & 0 \\
a'+c' & 1 
\end{pmatrix} - \sum_{b' \in \mathbb{F}_q}
\begin{pmatrix}
1 & b' \\
0 & 1 
\end{pmatrix}
- \sum_{a', c' \in \mathbb{F}_q^{\times}}
\begin{pmatrix}
-c'/a' & -(a'+c')/a'c' \\
0 & -a'/c' 
\end{pmatrix}
+\sum_{a' \in \mathbb{F}_q}
\begin{pmatrix}
1 & 0 \\
0 & 1 
\end{pmatrix}.\]

Simplifying similar terms and remembering that $\sum_{a \in \mathbb{F}_q} \begin{pmatrix}
1 & a \\
0 & 1 
\end{pmatrix} = qe_{U_{\lambda}}$, we get the following relation
\begin{multline*}
    q^3 e_{U_{\lambda}}e_{U_{\lambda}^{op}}e_{U_{\lambda}} = q^3 e_{U_{\lambda}^{op}}e_{U_{\lambda}}e_{U_{\lambda}^{op}} + (q-1) qe_{U_{\lambda}} - (q-1) qe_{U_{\lambda}^{op}} +\\
    +
    \sum_{a, c \in \mathbb{F}_q^{\times}}\begin{pmatrix}
    -a/c & 0 \\
    -(a+c)/ac & -c/a 
    \end{pmatrix}
    - \sum_{a', c' \in \mathbb{F}_q^{\times}}
    \begin{pmatrix}
    -c'/a' & -(a'+c')/a'c' \\
    0 & -a'/c' 
    \end{pmatrix}.
\end{multline*}

Now, we can write $\R^2$ as $ q^4 e_{U_{\lambda}}e_{U_{\lambda}^{op}}e_{U_{\lambda}}e_{U_{\lambda}}e_{U_{\lambda}^{op}}e_{U_{\lambda}}=q^4 e_{U_{\lambda}}e_{U_{\lambda}^{op}}e_{U_{\lambda}}e_{U_{\lambda}^{op}}e_{U_{\lambda}}$. Noting that $e_{U_{\lambda}}e_{U_{\lambda}^{op}}e_{U_{\lambda}}e_{U_{\lambda}^{op}}e_{U_{\lambda}} = e_{U_{\lambda}}(e_{U_{\lambda}^{op}}e_{U_{\lambda}}e_{U_{\lambda}^{op}})e_{U_{\lambda}}$, we can expand and get

\begin{multline*}
    \R^2 =  q^4e_{U_{\lambda}}e_{U_{\lambda}^{op}}e_{U_{\lambda}} -  (q-1) q^2e_{U_{\lambda}} +  (q-1) q^2e_{U_{\lambda}}e_{U_{\lambda}^{op}}e_{U_{\lambda}}-\\
    - qe_{U_{\lambda}} \sum_{a, c \in \mathbb{F}_q^{\times}}\begin{pmatrix}
    -a/c & 0 \\
    -(a+c)/ac & -c/a 
    \end{pmatrix} e_{U_{\lambda}}
    + qe_{U_{\lambda}} \sum_{a', c' \in \mathbb{F}_q^{\times}}
    \begin{pmatrix}
    -c'/a' & -(a'+c')/a'c' \\
    0 & -a'/c' 
    \end{pmatrix} e_{U_{\lambda}}.
\end{multline*}

Let's consider the first sum, the second sum is essentially identical. Fix a value of $a$ and $c$.\\

Now, consider the pair $ka, kc$. The matrix then becomes 
\begin{equation*}
    \begin{pmatrix}
-ka/kc & 0 \\
-(ka+kc)/kakc & -kc/ka 
\end{pmatrix} = \begin{pmatrix}
-a/c & 0 \\
-(a+c)/kac & -c/a 
\end{pmatrix}.    
\end{equation*}

Note that this is simply multiplying the bottom left element by $k^{-1}$. Thus, if $a+c \ne 0$, we can always make any value appear in the bottom left square. In addition, by varying $a$, we can obtain any value in the top left (other than $0$) with equal frequency. However, note that if $a+c = 0$, then the element on the diagonal is $1$. Thus, we can characterize all possible matrices in the following way: the element on the top left runs over all values in $\mathbb{F}_q^{\times}$ with equal frequency, the element in the bottom right is its inverse, and the element in the bottom left is in $\mathbb{F}_q^{\times}$ with equal frequency as long as $-a/c \ne 1$, and it is always $0$ otherwise. It is simple to check that all possible matrices show up exactly once with the exception of the identity matrix, which shows up $|\mathbb{F}_q^\times|$ times.\\

So we have
\[
\sum_{a, c \in \mathbb{F}_q^{\times}}\begin{pmatrix}
    -a/c & 0 \\
    -(a+c)/ac & -c/a 
    \end{pmatrix} = 
\sum_{\substack{a \in \mathbb{F}_q^\times \setminus \{1\}\\ t\in \mathbb{F}_q^{\times}}} \begin{pmatrix}
a & 0 \\
t & a^{-1}
\end{pmatrix} + (q-1) I.\]
%We can add and subtract the unipotent matrices (which have $1$s on the diagonal and $\mathbb{F}_q$ on the bottom left) to get

%\begin{align*}
%    \sum_{a, c \in \mathbb{F}_q^{\times}}\begin{pmatrix}
%    -a/c & 0 \\
%    -(a+c)/ac & -c/a 
%    \end{pmatrix} = \sum_{a \in \mathbb{F}_q^\times} \begin{pmatrix}
%    a & 0 \\
%    \mathbb{F}_q^\times & a^{-1}
%    \end{pmatrix} - V_\lambda^{op}+q I.
%\end{align*}
%Now, we can pull out the diagonal matrices to get
This simplifies to
\begin{align*}
    \sum_{a, c \in \mathbb{F}_q^{\times}}\begin{pmatrix}
    -a/c & 0 \\
    -(a+c)/ac & -c/a 
    \end{pmatrix} &= \sum_{a \in \mathbb{F}_q^\times} \begin{pmatrix}
    a & 0 \\
    0 & a^{-1}
    \end{pmatrix} (qe_{U_{\lambda}^{op}} - I)- qe_{U_{\lambda}^{op}}+q I.
\end{align*}
Plugging this back into our original equation we get

\begin{multline*}
    \R^2 =  q^4e_{U_{\lambda}}e_{U_{\lambda}^{op}}e_{U_{\lambda}} -  (q-1) q^2e_{U_{\lambda}} +  (q-1) q^2e_{U_{\lambda}}e_{U_{\lambda}^{op}}e_{U_{\lambda}}-\\
    - qe_{U_{\lambda}} \left( \sum_{a \in \mathbb{F}_q^\times} \begin{pmatrix}
    a & 0 \\
    0 & a^{-1}
    \end{pmatrix} (qe_{U_{\lambda}^{op}} - I)- qe_{U_{\lambda}^{op}}+q I\right) e_{U_{\lambda}}+\\
    + qe_{U_{\lambda}} \left(  \sum_{a \in \mathbb{F}_q^\times} \begin{pmatrix}
    a & 0 \\
    0 & a^{-1}
    \end{pmatrix} (qe_{U_{\lambda}} - I)- qe_{U_{\lambda}}+q I \right) e_{U_{\lambda}}.
\end{multline*}

We can distribute out and cancel a few terms:%, noting that $q = |U_\lambda|$:

\begin{multline*}
    \R^2 =  q^4e_{U_{\lambda}}e_{U_{\lambda}^{op}}e_{U_{\lambda}} - q^3 e_{U_{\lambda}} + q^3e_{U_{\lambda}}e_{U_{\lambda}^{op}}e_{U_{\lambda}} -\\- qe_{U_{\lambda}} \left( \sum_{a \in \mathbb{F}_q^\times} \begin{pmatrix}
    a & 0 \\
    0 & a^{-1}
    \end{pmatrix} (qe_{U_{\lambda}^{op}} - I)\right) e_{U_{\lambda}}+
     qe_{U_{\lambda}} \left(  \sum_{a \in \mathbb{F}_q^\times} \begin{pmatrix}
    a & 0 \\
    0 & a^{-1}
    \end{pmatrix} (qe_{U_{\lambda}} - I) \right) e_{U_{\lambda}}.
\end{multline*}

Now, we can factor out the Levi part from the equation

\begin{equation*}
     \R^2=  q^4e_{U_{\lambda}}e_{U_{\lambda}^{op}}e_{U_{\lambda}} - q^3 e_{U_{\lambda}} + q^3e_{U_{\lambda}}e_{U_{\lambda}^{op}}e_{U_{\lambda}} + q^2\left(  \sum_{a \in \mathbb{F}_q^\times} \begin{pmatrix}
    a & 0 \\
    0 & a^{-1}
    \end{pmatrix} \right) ( e_{U_{\lambda}} - e_{U_{\lambda}}e_{U_{\lambda}^{op}}e_{U_{\lambda}}).
\end{equation*}

Now, we can substitute back in $\R = q^2e_{U_{\lambda}}e_{U_{\lambda}^{op}}e_{U_{\lambda}}$ to get the next equality:
\begin{equation*}
    \R^2 =  q^2 \R - q^3 e_{U_{\lambda}} + q \R + \left(  \sum_{a \in \mathbb{F}_q^\times} \begin{pmatrix}
    a & 0 \\
    0 & a^{-1}
    \end{pmatrix} \right) ( q^2 e_{U_{\lambda}} - \R).
\end{equation*}

We move some terms to the other side
\begin{equation*}
    \R^2 -  q^2 \R + q^3 e_{U_{\lambda}} - q \R  = -\left(  \sum_{a \in \mathbb{F}_q^\times} \begin{pmatrix}
    a & 0 \\
    0 & a^{-1}
    \end{pmatrix} \right) ( \R - q^2 e_{U_{\lambda}})
\end{equation*}
and square both sides
\begin{multline*}
    \left(\R^2 -  q^2 \R + q^3 e_{U_{\lambda}} - q \R\right)^2 =\\
    \left(  \sum_{a \in \mathbb{F}_q^\times} \begin{pmatrix}
    a & 0 \\
    0 & a^{-1}
    \end{pmatrix} \right) (q^2 e_{U_{\lambda}} - \R) \left(  \sum_{b \in \mathbb{F}_q^\times} \begin{pmatrix}
    b & 0 \\
    0 & b^{-1}
    \end{pmatrix} \right) ( q^2 e_{U_{\lambda}} - \R).
\end{multline*}

However, since $\R$ and $e_{U_{\lambda}}$ commute with diagonal matrices, we can actually write it as follows:
\begin{equation*}
    \left(\R^2 -  q^2 \R + q^3 e_{U_{\lambda}} - q \R\right)^2 =\\
    \left(  \sum_{a \in \mathbb{F}_q^\times} \begin{pmatrix}
    a & 0 \\
    0 & a^{-1}
    \end{pmatrix} \right)^2 ( q^2 e_{U_{\lambda}} - \R)^2.
\end{equation*}

%Note that this set of Levi matrices form a group, so when we square it in our algebra, we simply get the order of the group multiplied by the group itself. In our case, it looks like
Therefore,
\begin{equation*}
    \left(\R^2 -  q^2 \R + q^3 e_{U_{\lambda}} - q \R\right)^2 =
    (q-1) \left(  \sum_{a \in \mathbb{F}_q^\times} \begin{pmatrix}
    a & 0 \\
    0 & a^{-1}
    \end{pmatrix} \right) ( q^2 e_{U_{\lambda}} - \R)^2.
\end{equation*}

Note that the right side has a multiple that looks very similar to the equation we had before we squared it, so we can substitute that back in, allowing us to get rid of the sum
\begin{equation*}
    \left(\R^2 -  q^2 \R + q^3 e_{U_{\lambda}} - q \R\right)^2 =
    (1-q) \left(\R^2 -  q^2 \R + q^3 e_{U_{\lambda}} - q \R\right) ( \R - q^2e_{U_{\lambda}} ),
\end{equation*}
\begin{equation*}
    \left(\R^2 -  q^2 \R + q^3 e_{U_{\lambda}} - q \R\right)
    \left(\R^2 -  q^2 \R + q^3 e_{U_{\lambda}} - q \R-(1-q)(\R - q^2e_{U_{\lambda}})\right) = 0.
\end{equation*}

Finally, since we are studying Radon transform on $U_\lambda$-invariant functions, we may substitute $1$ for $e_{U_{\lambda}}$, and we can simplify and factorize this equation.

\begin{equation*}
    (\R - q^2)^2(\R - q)(\R -1) = 0.
\end{equation*}
 
From this we get that all the possible eigenvalues of $\R$ are $1$, $q$, and $q^2$. An explicit computation (e.g. for prime $q$) shows that every eigenvalue actually appears in the spectrum.
\end{proof}

\end{document}
